\documentclass[10pt,a4paper]{amsart}
\usepackage[margin=19mm]{geometry}
\usepackage{amsmath,amssymb,mathtools}
\usepackage[scr=rsfs,cal=euler]{mathalfa}
\usepackage{xfrac}
\usepackage[unicode,hidelinks,bookmarks=false]{hyperref}
\newcommand{\ie}{i.e.\ }
\newcommand{\cf}{cf.\ }
\newcommand{\resp}{respectively\ }
\newcommand{\Subsection}{\subsection}
\providecommand{\nobreakdash}{\nobreak\hbox{-}}
\setlength{\emergencystretch}{2em}
\multlinegap0pt
\usepackage{booktabs}
\newtheorem{theorem}{Theorem}[section]
\newtheorem{lemma}[theorem]{Lemma}
\newtheorem{proposition}[theorem]{Proposition}
\newtheorem{definition}[theorem]{Definition}
\newtheorem{remark}[theorem]{Remark}
\newcommand{\card}[1]{\lvert #1\rvert}
\newcommand{\Rcal}{\mathcal{R}}
\newcommand{\Acal}{\mathcal{A}}
\begin{document}
\title[An Explicit Counterexample to Stanley's Rankwise Lower-Bound]{An Explicit Counterexample to Stanley's Rankwise Lower-Bound Conjecture for Differential Posets}
\author{Secondary 05B30; Xinan Dai, Wenhao Deng, Yingdong Shi, Tailin Wu, Yuchen Yang}
\date{}
\maketitle
\noindent\emph{Source and licence.} An Explicit Counterexample to Stanley's Rankwise Lower-Bound Conjecture for Differential Posets. Version arXiv:2607.22988v2 (CC BY 4.0). This material is distributed under the Creative Commons Attribution 4.0 International licence, http://creativecommons.org/licenses/by/4.0/, as recorded in the arXiv OAI licence metadata for this exact version and shown on the abstract page https://arxiv.org/abs/2607.22988v2. Full rights notice: licenses/arxiv-2607.22988-CC-BY-4.0.txt.\par\smallskip
\noindent\textit{Editorial scope.} Counted unit: the original \texttt{lemma} environment \texttt{lem:point-pair} at source lines 168--177 together with its complete proof at lines 179--181; the delivered document keeps source lines 92--306 so that every referenced label and every citation resolves inside it. Native editable source is the paper's own concatenated TeX; the AMS wrapper is editorial formatting only. No publisher class, upstream script or unrelated artwork is redistributed.
\begin{equation}\label{eq:stanley-bound}
  \card{P_n}\geq \card{(Y^r)_n}
  \qquad(n\geq0)?
\end{equation}

We show that the universally quantified assertion in \eqref{eq:stanley-bound} is false. The basic counterexample occurs at $r=3$ and rank four. Its finite mechanism is an unequal replacement of lower-cover blocks: thirteen blocks in $Y^3$ are replaced by twelve new blocks while preserving every one-point and two-point incidence multiplicity. These multiplicities are precisely the entries of the operator $D_4U_3$. Consequently, the differential-poset relation remains valid through rank three, although the fourth rank has one fewer element. Stanley's reflection construction then extends the finite initial segment to an infinite differential poset.

Section~\ref{sec:preliminaries} recalls the relevant definitions and states the main result. Section~\ref{sec:r3-construction} gives the rank-four replacement for $r=3$ and verifies it directly. Section~\ref{sec:extension} gives the infinite extension, and Section~\ref{sec:general-r} transports the construction to every $r\geq3$.

\section{Differential posets and the main result}\label{sec:preliminaries}

Young's lattice $Y$ is the graded poset of integer partitions ordered by inclusion of Ferrers diagrams. Its rank-$n$ elements are the partitions of $n$, and a cover relation adds one cell. Write $p(n)=\card{Y_n}$ for the number of integer partitions of $n$. Hence
\[
  F_Y(q)=\sum_{n\geq0}\card{Y_n}q^n
  =\prod_{i\geq1}(1-q^i)^{-1}.
\]
The Cartesian product of $r$ copies is denoted by $Y^r$. An element of $(Y^r)_n$ is an $r$-tuple of partitions whose total size is $n$, and
\[
  F_{Y^r}(q)=F_Y(q)^r.
\]

\begin{definition}\label{def:differential-poset}
Let $r$ be a positive integer. An \emph{$r$-differential poset} is a locally finite graded poset
\[
  P=\bigsqcup_{n\geq0}P_n
\]
with a unique least element and the following properties.
\begin{enumerate}
  \item[\textup{(D1)}] If $x\neq y$ have the same rank, then the number of their common upper covers equals the number of their common lower covers.
  \item[\textup{(D2)}] If $x$ has $d$ lower covers, then it has $d+r$ upper covers.
\end{enumerate}
\end{definition}

Let $\mathbb{Q}P_n$ be the vector space with basis $P_n$. Define the up and down operators by
\[
  U_nx=\sum_{\substack{z\in P_{n+1}\\z\gtrdot x}}z,
  \qquad
  D_nx=\sum_{\substack{w\in P_{n-1}\\x\gtrdot w}}w.
\]
The two axioms are equivalent to
\begin{equation}\label{eq:differential-identity}
  D_{n+1}U_n-U_{n-1}D_n=rI_{\mathbb{Q}P_n}
  \qquad(n\geq0),
\end{equation}
with the evident convention in rank zero. Young's lattice is $1$-differential, and Cartesian products add the parameters; hence $Y^r$ is $r$-differential.

\begin{theorem}\label{thm:main}
For every integer $r\geq3$, there exists an infinite $r$-differential poset $P^{(r)}$ such that
\[
  \card{P^{(r)}_i}=\card{(Y^r)_i}\quad(0\leq i\leq3),
  \qquad
  \card{P^{(r)}_4}=\card{(Y^r)_4}-\left\lfloor\frac r3\right\rfloor.
\]
Moreover,
\begin{equation}\label{eq:Y-r-rank4}
  \card{(Y^r)_4}=\frac{r(r+1)(r^2+17r+42)}{24}.
\end{equation}
In particular, the first five rank sizes of $P^{(3)}$ are
\[
  (1,3,9,22,50),
\]
whereas those of $Y^3$ are $(1,3,9,22,51)$.
\end{theorem}

The rank-growth results of Miller~\cite{Miller2013} and Gaetz--Venkataramana~\cite{GaetzVenkataramana2020} do not imply the coefficientwise comparison in \eqref{eq:stanley-bound} and are consistent with Theorem~\ref{thm:main}.

\section{The finite \texorpdfstring{$r=3$}{r=3} construction}\label{sec:r3-construction}

\subsection{Lower-cover blocks}

Retain ranks zero through three of $Y^3$. For a proposed rank-four element $z$, set
\[
  B_z=\{x\in(Y^3)_3:x\lessdot z\}.
\]
We call $B_z$ the \emph{lower-cover block} of $z$.


\begin{lemma}[Point-and-pair criterion]\label{lem:point-pair}
Fix the poset through rank three. Let $(B_z)_{z\in Q}$ be a finite indexed family of nonempty subsets of $P_3$, one for each proposed rank-four element $z\in Q$. For $x,y\in P_3$, the $(x,y)$-entry of $D_4U_3$ is
\[
  \begin{cases}
    \card{\{z\in Q:x\in B_z\}},&x=y,\\[2mm]
    \card{\{z\in Q:\{x,y\}\subseteq B_z\}},&x\neq y.
  \end{cases}
\]
Consequently, replacing rank-four blocks by blocks having the same incidence multiplicity on every one- and two-element subset of $P_3$ leaves $D_4U_3$ unchanged.
\end{lemma}

\begin{proof}
For $x=y$, the coefficient of $x$ in $D_4U_3x$ is the number of rank-four elements covering $x$. For $x\neq y$, the coefficient of $y$ is the number of rank-four elements covering both $x$ and $y$.
\end{proof}


Because the ranks below four are unchanged, $U_2D_3$ is unchanged as well. It therefore suffices to preserve the point and pair incidence multiplicities of the rank-four lower-cover blocks.

\subsection{Indexing rank three}

Write a triple of partitions as $(\lambda;\mu;\nu)$, write $0$ for the empty partition, and concatenate the parts of a partition; for example, $21$ denotes $(2,1)$. We label the elements of $(Y^3)_3$ as follows.

\begin{center}
\renewcommand{\arraystretch}{1.08}
\begin{tabular}{@{}r l r l@{}}
\toprule
Index & Element & Index & Element\\
\midrule
0  & $(3;0;0)$   & 11 & $(1;0;11)$\\
1  & $(21;0;0)$  & 12 & $(0;3;0)$\\
2  & $(111;0;0)$ & 13 & $(0;21;0)$\\
3  & $(2;1;0)$   & 14 & $(0;111;0)$\\
4  & $(2;0;1)$   & 15 & $(0;2;1)$\\
5  & $(11;1;0)$  & 16 & $(0;11;1)$\\
6  & $(11;0;1)$  & 17 & $(0;1;2)$\\
7  & $(1;2;0)$   & 18 & $(0;1;11)$\\
8  & $(1;11;0)$  & 19 & $(0;0;3)$\\
9  & $(1;1;1)$   & 20 & $(0;0;21)$\\
10 & $(1;0;2)$   & 21 & $(0;0;111)$\\
\bottomrule
\end{tabular}
\end{center}

\subsection{The thirteen-for-twelve replacement}\label{subsec:replacement}

Begin with the $51$ lower-cover blocks supplied by rank four of $Y^3$. Delete
\begin{equation}\label{eq:R-blocks}
\begin{aligned}
\Rcal=\{&\{1\},\{13\},\{20\},\{5,7\},\{5,8\},\{6,10\},\{6,11\},\\
&\{1,3,5\},\{1,4,6\},\{3,4,9\},\{5,6,9\},\{7,8,13\},\{10,11,20\}\}.
\end{aligned}
\end{equation}
In the displayed order, these are the lower-cover blocks of
\begin{equation}\label{eq:R-elements}
\begin{aligned}
 &(22;0;0),\ (0;22;0),\ (0;0;22),\ (11;2;0),\ (11;11;0),\ (11;0;2),\ (11;0;11),\\
 &(21;1;0),\ (21;0;1),\ (2;1;1),\ (11;1;1),\ (1;21;0),\ (1;0;21).
\end{aligned}
\end{equation}
The list in \eqref{eq:R-elements} is checked directly by deleting one removable corner cell from one of the three partition coordinates.
Insert twelve new rank-four elements with lower-cover blocks
\begin{equation}\label{eq:A-blocks}
\begin{aligned}
\Acal=\{&\{1,5\},\{1,6\},\{5,6\},\{7,13\},\{8,13\},\{10,20\},\{11,20\},\\
&\{1,3,4\},\{3,5,9\},\{4,6,9\},\{5,7,8\},\{6,10,11\}\}.
\end{aligned}
\end{equation}
Every block in \eqref{eq:A-blocks} is nonempty and therefore defines a rank-four element once the indicated cover relations are declared.

\subsection{Verification}

Only the twelve vertices
\[
  1,3,4,5,6,7,8,9,10,11,13,20
\]
occur in $\Rcal\cup\Acal$. Their incidence multiplicities agree:
\begin{equation}\label{eq:point-incidences}
\begin{array}{c|rrrrrrrrrrrr}
 v&1&3&4&5&6&7&8&9&10&11&13&20\\
\hline
 \deg_{\Rcal}(v)=\deg_{\Acal}(v)&3&2&2&4&4&2&2&2&2&2&2&2.
\end{array}
\end{equation}
Both sides contain each of the following twenty-two pairs exactly once:
\begin{equation}\label{eq:pair-incidences}
\begin{aligned}
&\{1,3\},\{1,4\},\{1,5\},\{1,6\},\{3,4\},\{3,5\},\{3,9\},\{4,6\},\{4,9\},\{5,6\},\\
&\{5,7\},\{5,8\},\{5,9\},\{6,9\},\{6,10\},\{6,11\},\{7,8\},\{7,13\},\{8,13\},\\
&\{10,11\},\{10,20\},\{11,20\},
\end{aligned}
\end{equation}
and neither side contains any other pair. Equations~\eqref{eq:point-incidences} and~\eqref{eq:pair-incidences} are therefore a complete incidence check.

Let $P^{(3)}_{\leq4}$ be obtained by leaving $(Y^3)_{\leq3}$ unchanged and replacing $\Rcal$ by $\Acal$ in rank four. By Lemma~\ref{lem:point-pair}, its $D_4U_3$ matrix is the same as that of $Y^3$, while $U_2D_3$ is unchanged. Hence
\[
  D_4U_3-U_2D_3=3I
\]
on rank three. The lower-rank identities are inherited from $Y^3$, so $P^{(3)}_{\leq4}$ is a partial $3$-differential poset and
\[
  \card{P^{(3)}_4}=51-13+12=50.
\]

\begin{remark}\label{rem:trade}
The replacement has a design-theoretic interpretation. Define a signed function on subsets of the twelve affected vertices by
\[
  f(B)=
  \begin{cases}
    +1,&B\in\Rcal,\\
    -1,&B\in\Acal,\\
    0,&\text{otherwise}.
  \end{cases}
\]
Then
\[
  \sum_{B\supseteq S}f(B)=0
  \qquad(1\leq\card{S}\leq2),
\]
whereas $\sum_Bf(B)=1$. After adjoining the empty block to the $\Acal$-side, the two collections form a simple $[2]$-trade of volume $13$ in the sense of~\cite{GhorbaniEtAl2020}. The differential-poset axioms detect the first two incidence moments but not the zeroth, which accounts for the one-element saving.
\end{remark}

\section{Extension to an infinite poset}\label{sec:extension}

The construction above is finite. The following reflection extension, appearing in~\cite[Proposition~6.1]{Stanley1988}, preserves its initial ranks.

\begin{proposition}[Reflection extension]\label{prop:reflection-extension}
Suppose $P_{\leq n}$ is a finite partial $r$-differential poset of rank $n$, meaning that \eqref{eq:differential-identity} holds in ranks below $n$. Then $P_{\leq n}$ extends by one rank to a partial $r$-differential poset of rank $n+1$. Iteration produces an infinite $r$-differential poset whose first $n+1$ ranks, including their cover relations, are exactly $P_{\leq n}$.
\end{proposition}

\begin{proof}
For each $y\in P_{n-1}$, add an element $y^*$ of rank $n+1$ that covers precisely the elements $x\in P_n$ covering $y$. For each $x\in P_n$, also add $r$ rank-$(n+1)$ elements, each covering only $x$.

If $x\in P_n$ has $d$ lower covers, then the elements $y^*$ provide $d$ upper covers and the singleton elements provide $r$ more. Thus $x$ has $d+r$ upper covers. If $x\neq x'$ lie in $P_n$, their common new upper covers are precisely the elements $y^*$ indexed by their common lower covers. The differential conditions therefore hold at rank $n$, and all lower-rank relations are unchanged. Iterating and taking the union gives the required infinite poset.
\end{proof}

Applying Proposition~\ref{prop:reflection-extension} to $P^{(3)}_{\leq4}$ gives an infinite $3$-differential poset with initial rank sequence
\[
  1,3,9,22,50.
\]

\section{The construction for every \texorpdfstring{$r\geq3$}{r >= 3}}\label{sec:general-r}

\begin{thebibliography}{99}
\bibitem[GaetzVenkataramana2020]{GaetzVenkataramana2020} C.~Gaetz and P.~Venkataramana, \emph{Path counting and rank gaps in differential posets}, Order \textbf{37} (2020), no.~2, 279--286, arXiv:1806.03509.
\bibitem[GhorbaniEtAl2020]{GhorbaniEtAl2020} E.~Ghorbani, S.~Kamali, G.~B.~Khosrovshahi, and D.~S.~Krotov, \emph{On the volumes and affine types of trades}, Electron. J. Combin. \textbf{27} (2020), no.~1, Paper No.~P1.29, arXiv:1810.02296.
\bibitem[Miller2013]{Miller2013} A.~R.~Miller, \emph{Differential posets have strict rank growth: a conjecture of Stanley}, Order \textbf{30} (2013), no.~2, 657--662, arXiv:1202.3006.
\bibitem[Stanley1988]{Stanley1988} R.~P.~Stanley, \emph{Differential posets}, J. Amer. Math. Soc. \textbf{1} (1988), no.~4, 919--961.
\end{thebibliography}
\end{document}
